\section{Related Work}
\label{sec:related-work}
% Source status: the two author-supplied Lim manuscripts are cited as
% manuscripts, without inferring a publication venue or an unknown year.
% This section describes planned positioning, not completed experiments.

\subsection{Traffic Complexity and Operational Validation}
\label{sec:rw-complexity}
Historical trajectories provide observable evidence of how arrival traffic
has been managed. Lim et al.~\cite{lim2024datadriven} quantify terminal manoeuvring
area (TMA) complexity through vectoring and holding, using normalized arc
length, normalized total curvature, and holding duration. Their framework
applies principal component analysis (PCA) to these three features,
clusters trajectories with $k$-means, ranks the clusters by operational
complexity, and aggregates their assigned weights over one-hour scenes.
This construction incorporates both traffic volume and the prevalence of
operational interventions. It provides an interpretable reference for
scene assessment, although its feature set explicitly determines which
aspects of trajectory behaviour contribute to the score.

The subsequent retrospective study by Lim et al.~\cite{lim2026retrospective}
compares the model with six rated approach controllers assessing twenty
historical one-hour scenarios, with complementary eye-tracking measures.
The continuous model score achieves a mean Spearman correlation of
$0.746$ with individual controller ratings. These results support its
operational relevance within the studied TMA, while the retrospective
playback setting and limited scenario set constrain broader conclusions.
Together, these studies establish useful operational comparators for a
new geometric descriptor. They also motivate a distinction between
trajectory-derived complexity and controller workload: agreement with
trajectory features alone does not establish agreement with human
judgements or workload under active control.

\subsection{Self-Supervised Aircraft Trajectory Representations}
\label{sec:rw-representation}
Aircraft Trajectory Segmentation-Based Contrastive Coding (ATSCC)
learns representations from unlabelled trajectories using a causal
Transformer encoder \cite{phisannupawong2025atscc}. It derives segment
identities through Ramer--Douglas--Peucker simplification and uses a
modified soft nearest-neighbour objective to bring representations from
the same segment together while separating other segments and
trajectories. Its route classification and clustering evaluations show
that this supervision can recover meaningful trajectory structure.
However, route discrimination does not by itself validate a
representation of operational complexity.

Lim et al.~\cite{limcontrastive} directly investigate this distinction by probing
ATSCC-derived embeddings for holding detection and regression of
trajectory-based operational indicators. Their comparison identifies
temporal pooling as a major determinant of downstream performance, with
sum pooling performing particularly well on complexity-related regression
tasks. Crucially, evaluation within individual standard terminal arrival
routes (STARs) reveals weaker sensitivity to some operational variations:
holding-duration $R^2$ falls from $0.6820$ in the global evaluation to
$0.3315$ within STARs, whereas holding-detection AUC remains $0.9539$.
This evidence motivates evaluation beyond route identity and scrutiny of
duration-related information introduced by pooling. For the present
study, an additional question is whether representations learned from
trajectories retain useful behavioural differences when applied to
fragments clipped to hourly windows. The existing evaluations do not
establish the reliability of a geometric statistic computed across those
fragment embeddings.

\subsection{Correlation Dimension and Learned Point-Cloud Geometry}
\label{sec:rw-dimension}
Grassberger and Procaccia~\cite{grassberger1983measuring} introduced the correlation-dimension
approach, which characterizes how the probability of two points lying
within distance $r$ scales with $r$. In empirical applications, the slope
of a log--log correlation-sum curve over a suitable scale interval
provides an estimate of this scaling behaviour. This offers a way to
describe the geometry of an embedding population beyond its number of
coordinates. Related work on neural representations also demonstrates
that estimated intrinsic dimension can differ substantially from layer
width and change across learned representations
\cite{ansuini2019intrinsic}. Such findings motivate studying embedding
geometry, but do not imply that its dimension measures operational
difficulty in air traffic.

Finite-sample estimation is a central limitation.
Erba et al.~\cite{erba2019intrinsic} show how local undersampling can bias
correlation-based dimension estimates, particularly as intrinsic
dimension increases, and propose an estimator based on the full
correlation integral under additional geometric assumptions. A
complementary approach, TWO-NN, estimates intrinsic dimension from ratios
of first- and second-nearest-neighbour distances
\cite{facco2017estimating}. These approaches provide methodological
comparisons rather than interchangeable definitions: their assumptions
and finite-sample behaviour differ, and agreement is not guaranteed for
heterogeneous traffic scenes. In particular, reducing the encoder output
dimension does not establish that an hourly point cloud contains enough
aircraft for reliable estimation. We therefore treat the fitted
correlation slope as a finite-scale empirical descriptor whose
stability must be assessed against sample size, scale selection, and
representation choices.

\subsection{Positioning of the Present Study}
\label{sec:rw-positioning}
The proposed unit of analysis is an hourly scene represented by a point
cloud with one embedding per aircraft trajectory fragment. The points
describe a population of aircraft behaviours; they are not delay
coordinates reconstructed from a single dynamical trajectory, and the
analysis makes no claim about chaotic traffic dynamics. We investigate
whether empirical correlation dimension captures scene variation that
adds information beyond aircraft count, route composition, and simple
embedding dispersion. A fixed encoder, pooling rule, and distance
definition make comparisons across scenes interpretable. The proposed
validation separates geometric stability from operational relevance,
using traffic-volume and route-composition controls, fragment-duration
and coverage checks, and independent operational or controller-based
comparisons where available. A higher estimated dimension is thus a
hypothesis about richer latent structure, whose relationship to
operational complexity remains to be established.
